\section{Where the Divergence Is Built: The First Decade of the Career}\label{sec:careers}

If the divergence is a divergence in wage growth, it can arise between jobs---sons from richer families moving into
better-paid occupations---or within them, through faster wage growth in the job a son holds. This section follows the
order in which the evidence answers the question: how much is occupational upgrading, how much is predicted by the job
a son holds at 25--34 and who gets those jobs, how wages grow within them, and what role schooling plays. All results are
descriptive: sons select into schooling, occupations and employers, and nothing here identifies a causal effect of
parental income or race. Racial differences are White minus Black.

\subsection{About one fifth is occupational upgrading}

All PSID occupation codes are mapped to David Dorn's consistent occ1990dd scheme, and each occupation gets a wage score:
the mean log hourly wage of men aged 25--64 in the occupation relative to all men, from 2.9 million CPS Outgoing
Rotation Group records of 1979--2019. Each son's log wage is the sum of the score of his occupation and his wage
relative to it, so the parental-income gradient of the log wage splits exactly into a between- and a
within-occupation part.\footnote{Code: \texttt{IGE/27.CareerWageGrowth.py} (\texttt{output/career\_wage\_growth.md},
Section 0). The score accounts for about a quarter of the variance of sons' log wages.} Pooled, the gradient rises
from 0.259$^{***}$ at 25--34 to 0.410$^{***}$ at 35--55; the between-occupation part barely moves (0.158 to 0.176) and
the within-occupation part more than doubles (0.102 to 0.235) (Figure~\ref{fig:career}, Table~\ref{tab:careerdecomp}).
By race, the occupational gap is in place at 25--34: the White--Black difference in the wage gradient, 0.090$^{***}$, is
almost all between occupations (0.082$^{***}$). From 25--34 to 35--55 the difference grows by 0.167$^{***}$ (0.037), of
which 0.031$^{*}$ is the occupational score (a share of 0.18$^{**}$; 0.27$^{***}$ with a time-varying score) and
0.136$^{***}$ is within occupations. Background buys a better-paid occupation at the start, and the occupational
advantage hardly grows afterwards; what grows is the wage within the occupation.

\begin{figure}[!htbp]
\centering
\includegraphics[width=\textwidth]{career_wage_growth.png}
\caption{Wage persistence by the son's age: between and within occupations (PSID)}\label{fig:career}
\begin{minipage}{\textwidth}\footnotesize \textit{Notes:} Coefficient of the son's log wage, of his occupation's wage
score and of his within-occupation residual on log parental income, by age band, with 95\% confidence intervals. The
score and the residual add up to the log wage. Code: \texttt{IGE/27.CareerWageGrowth.py}.\end{minipage}
\end{figure}

\begin{table}[!htbp]
\centering
\caption{Wage persistence by age: between and within occupations (PSID, sons)}
\label{tab:careerdecomp}
\resizebox{\ifdim\width>\textwidth\textwidth\else\width\fi}{!}{%
\resizebox*{!}{\ifdim\totalheight>0.85\textheight 0.85\textheight\else\totalheight\fi}{%
\begin{tabular}{lcccc}
\toprule
Panel / row & \begin{tabular}[b]{@{}c@{}}ln wage\end{tabular} & \begin{tabular}[b]{@{}c@{}}Occupation\\score\end{tabular} & \begin{tabular}[b]{@{}c@{}}Within-\\occupation\\residual\end{tabular} & \begin{tabular}[b]{@{}c@{}}Share\\between\\(2)/(1)\end{tabular} \\
 & (1) & (2) & (3) & (4) \\
\midrule
\multicolumn{5}{l}{\textit{Pooled (weighted): coefficient on log parental income}} \\
Sons aged 25--34 & $0.259^{***}$ & $0.158^{***}$ & $0.102^{***}$ & $0.608^{***}$ \\
 & (0.017) & (0.010) & (0.014) & (0.038) \\
Sons aged 35--44 & $0.401^{***}$ & $0.174^{***}$ & $0.227^{***}$ & $0.433^{***}$ \\
 & (0.026) & (0.012) & (0.021) & (0.027) \\
Sons aged 45--55 & $0.429^{***}$ & $0.179^{***}$ & $0.250^{***}$ & $0.416^{***}$ \\
 & (0.039) & (0.017) & (0.033) & (0.037) \\
Sons aged 35--55 & $0.410^{***}$ & $0.176^{***}$ & $0.235^{***}$ & $0.428^{***}$ \\
 & (0.028) & (0.013) & (0.022) & (0.026) \\
Son-years (25--34 / 35--44 / 45--55) & 21,315 / 12,648 / 5,728 &  &  &  \\
\addlinespace
\multicolumn{5}{l}{\textit{Black: coefficient on log parental income}} \\
Sons aged 25--34 & $0.152^{***}$ & $0.073^{***}$ & $0.079^{***}$ & $0.480^{***}$ \\
 & (0.020) & (0.011) & (0.017) & (0.064) \\
Sons aged 35--44 & $0.159^{***}$ & $0.056^{***}$ & $0.103^{***}$ & $0.351^{***}$ \\
 & (0.029) & (0.014) & (0.023) & (0.069) \\
Sons aged 45--55 & $0.160^{***}$ & $0.067^{***}$ & $0.093^{**}$ & $0.420^{***}$ \\
 & (0.046) & (0.020) & (0.037) & (0.111) \\
Sons aged 35--55 & $0.161^{***}$ & $0.060^{***}$ & $0.101^{***}$ & $0.371^{***}$ \\
 & (0.030) & (0.014) & (0.024) & (0.068) \\
Son-years (25--34 / 35--44 / 45--55) & 7,412 / 4,149 / 1,843 &  &  &  \\
\addlinespace
\multicolumn{5}{l}{\textit{White: coefficient on log parental income}} \\
Sons aged 25--34 & $0.243^{***}$ & $0.155^{***}$ & $0.088^{***}$ & $0.638^{***}$ \\
 & (0.020) & (0.012) & (0.017) & (0.049) \\
Sons aged 35--44 & $0.405^{***}$ & $0.171^{***}$ & $0.234^{***}$ & $0.423^{***}$ \\
 & (0.032) & (0.014) & (0.026) & (0.030) \\
Sons aged 45--55 & $0.456^{***}$ & $0.175^{***}$ & $0.281^{***}$ & $0.383^{***}$ \\
 & (0.047) & (0.022) & (0.038) & (0.039) \\
Sons aged 35--55 & $0.418^{***}$ & $0.172^{***}$ & $0.246^{***}$ & $0.411^{***}$ \\
 & (0.033) & (0.015) & (0.026) & (0.028) \\
Son-years (25--34 / 35--44 / 45--55) & 13,915 / 8,500 / 3,885 &  &  &  \\
\addlinespace
\multicolumn{5}{l}{\textit{White $-$ Black: coefficient on log parental income}} \\
Sons aged 25--34 & $0.090^{***}$ & $0.082^{***}$ & $0.009$ & $0.903^{***}$ \\
 & (0.028) & (0.016) & (0.024) & (0.238) \\
Sons aged 35--44 & $0.246^{***}$ & $0.115^{***}$ & $0.131^{***}$ & $0.469^{***}$ \\
 & (0.043) & (0.020) & (0.035) & (0.071) \\
Sons aged 45--55 & $0.296^{***}$ & $0.108^{***}$ & $0.188^{***}$ & $0.364^{***}$ \\
 & (0.066) & (0.030) & (0.053) & (0.083) \\
Sons aged 35--55 & $0.258^{***}$ & $0.112^{***}$ & $0.145^{***}$ & $0.436^{***}$ \\
 & (0.045) & (0.020) & (0.035) & (0.065) \\
\bottomrule
\end{tabular}}}
\par\smallskip
\begin{minipage}{\textwidth}\footnotesize \textit{Continued on the next table; notes below its last part.}\end{minipage}
\end{table}
\begin{table}[!htbp]
\centering
\addtocounter{table}{-1}
\caption{Wage persistence by age: between and within occupations (PSID, sons) (continued)}
\resizebox{\ifdim\width>\textwidth\textwidth\else\width\fi}{!}{%
\resizebox*{!}{\ifdim\totalheight>0.85\textheight 0.85\textheight\else\totalheight\fi}{%
\begin{tabular}{lcccc}
\toprule
Panel / row & \begin{tabular}[b]{@{}c@{}}ln wage\end{tabular} & \begin{tabular}[b]{@{}c@{}}Occupation\\score\end{tabular} & \begin{tabular}[b]{@{}c@{}}Within-\\occupation\\residual\end{tabular} & \begin{tabular}[b]{@{}c@{}}Share\\between\\(2)/(1)\end{tabular} \\
 & (1) & (2) & (3) & (4) \\
\midrule
\multicolumn{5}{l}{\textit{Change in the coefficient across age bands (career divergence)}} \\
Pooled (weighted): (35--55) $-$ (25--34) & $0.151^{***}$ & $0.018^{*}$ & $0.133^{***}$ & $0.119^{*}$ \\
 & (0.023) & (0.010) & (0.021) & (0.062) \\
Pooled (weighted): (45--55) $-$ (25--34) & $0.169^{***}$ & $0.021$ & $0.148^{***}$ & $0.124$ \\
 & (0.035) & (0.015) & (0.032) & (0.081) \\
Black: (35--55) $-$ (25--34) & $0.009$ & $-0.013$ & $0.022$ &  \\
 & (0.025) & (0.011) & (0.022) &  \\
Black: (45--55) $-$ (25--34) & $0.008$ & $-0.006$ & $0.014$ &  \\
 & (0.042) & (0.019) & (0.036) &  \\
White: (35--55) $-$ (25--34) & $0.176^{***}$ & $0.017$ & $0.158^{***}$ & $0.099$ \\
 & (0.027) & (0.012) & (0.024) & (0.063) \\
White: (45--55) $-$ (25--34) & $0.213^{***}$ & $0.020$ & $0.193^{***}$ & $0.095$ \\
 & (0.042) & (0.020) & (0.037) & (0.084) \\
White $-$ Black: (35--55) $-$ (25--34) & $0.167^{***}$ & $0.031^{*}$ & $0.136^{***}$ & $0.184^{**}$ \\
 & (0.037) & (0.016) & (0.032) & (0.087) \\
White $-$ Black: (45--55) $-$ (25--34) & $0.206^{***}$ & $0.026$ & $0.179^{***}$ & $0.127$ \\
 & (0.060) & (0.027) & (0.051) & (0.117) \\
\bottomrule
\end{tabular}}}
\par\smallskip
\begin{minipage}{\textwidth}\footnotesize \textit{Notes:} One regression per cell: the son's outcome (column header) on log parental income (mean at 13--19), son-years in the age band, controls for quartics in the father's age at 16 and the son's age. Occupational wage score: mean log hourly wage of men aged 25--64 in the occupation relative to all men (CPS MORG, base 1992--2002), occupations harmonised to Dorn's occ1990dd scheme; log wage = score (between occupations) + residual (within occupation), so the score and residual columns add up to the log-wage column on the same sample. With a time-varying score (the same relative premium in $t-1$ to $t+1$) the between share of the change of the racial gap is 0.27 (\texttt{output/career\_wage\_growth.md}). Column (4) is left blank when the denominator is not different from zero. All cells are estimated as one stacked system, so differences across bands and groups have standard errors (clustered by 1968 family). Pooled: Black and White sons, weighted; Black and White: unweighted. Sons select into occupations, so the split is accounting, not causal. $^{***}$ $p<0.01$, $^{**}$ $p<0.05$, $^{*}$ $p<0.1$. Source: \texttt{27.CareerWageGrowth.py} (\texttt{output/career\_wage\_growth.csv}); table built by \texttt{28.PaperBTables.py}.\end{minipage}
\end{table}


\subsection{Where sons start: the job at 25--34}\label{sec:earlystart}

A within-occupation divergence can still be produced by where a son starts, if the jobs that sons of richer White
parents enter pay off later. I classify each son by the job he held at 25--34---his modal cell of 11 occupation groups
$\times$ 33 industry groups---whatever he does later, and add fixed effects for that early job to the regression of his
log wage on parental income, separately in each age band. The early job may then have any wage path afterwards, and
the drop in the parental-income coefficient it produces at 35--55 relative to 25--34 is the part of the rise of
persistence predicted by where the son started (Table~\ref{tab:earlyjob}).\footnote{93 percent of sons observed at
35--55 have an early job, and they look like all sons at 25--34. Randomly reassigning cells across sons gives shares of
about 0.03. Code: \texttt{IGE/35.EarlyIndustryLatePayoff.py} (\texttt{output/early\_industry\_late\_payoff.md}).}

The son's occupation-by-industry cell at 25--34 accounts for 0.405$^{***}$ of the pooled rise of persistence after 35
and 0.493$^{***}$ ($p=0.004$) of the racial divergence, mostly through occupation (his occupation group alone: 0.377$^{***}$
and 0.318$^{***}$); his detailed occupation accounts for 0.591$^{***}$ and 0.713$^{***}$. His industry group alone accounts
for 0.135$^{*}$ of the pooled rise and 0.108 ($p=0.30$) of the racial divergence, and the career slopes of industries,
estimated in the CPS, for 0.05--0.13 of either.\footnote{Appendix~\ref{app:careers}, Tables~\ref{tab:industrydecomp}
and \ref{tab:industrynonlinear}. Finance, insurance and legal services have relative wages that rise steeply from
25--29 to 35--39 and then flatten; their pay-off comes in the early thirties.} The early job also reflects the son's
own skills, so these shares measure how much where a son stands at 25--34 predicts how far he pulls away after 35, not a
causal return to the job. Within sons, the extra wage growth of well-paid starting jobs comes before 35: relative to
sons who started in cells that are neither well paid nor fast growing, their wages grow faster by 0.33--0.39 log points
per ten years between 25 and 35, and by 0.012--0.059 after 35.

\begin{table}[!htbp]
\centering
\caption{Where sons start: the job at 25--34 and the rise of wage persistence after 35 (PSID, sons)}
\label{tab:earlyjob}
\resizebox{\ifdim\width>\textwidth\textwidth\else\width\fi}{!}{%
\resizebox*{!}{\ifdim\totalheight>0.85\textheight 0.85\textheight\else\totalheight\fi}{%
\begin{tabular}{lcccccc}
\toprule
 & \begin{tabular}[b]{@{}c@{}}Accounted\\for at\\25--34\end{tabular} & \begin{tabular}[b]{@{}c@{}}Accounted\\for at\\35--55\end{tabular} & \begin{tabular}[b]{@{}c@{}}Change\end{tabular} & \begin{tabular}[b]{@{}c@{}}Share of\\pooled rise\end{tabular} & \begin{tabular}[b]{@{}c@{}}Change,\\White $-$ Black\end{tabular} & \begin{tabular}[b]{@{}c@{}}Share of\\gap change\end{tabular} \\
 & (1) & (2) & (3) & (4) & (5) & (6) \\
\midrule
\multicolumn{7}{l}{\textit{Fixed effects for the son's job at 25--34, by age band}} \\
Finance, insurance, legal, professional (example; 2 groups) & $0.014^{***}$ & $0.027^{***}$ & $0.013^{**}$ & $0.097^{**}$ & $0.012^{*}$ & $0.085^{*}$ \\
 & (0.003) & (0.007) & (0.005) & (0.039) & (0.006) & (0.046) \\
Industry groups ranked by pay and later growth (4 sets) & $0.024^{***}$ & $0.029^{***}$ & $0.005$ & $0.034$ & $0.001$ & $0.006$ \\
 & (0.005) & (0.008) & (0.005) & (0.037) & (0.008) & (0.058) \\
Industry group (33) & $0.034^{***}$ & $0.052^{***}$ & $0.018^{*}$ & $0.135^{*}$ & $0.015$ & $0.108$ \\
 & (0.007) & (0.013) & (0.010) & (0.072) & (0.016) & (0.105) \\
Manager or professional (2 groups) & $0.042^{***}$ & $0.076^{***}$ & $0.035^{***}$ & $0.257^{***}$ & $0.035^{***}$ & $0.246^{***}$ \\
 & (0.005) & (0.010) & (0.008) & (0.064) & (0.010) & (0.086) \\
Occupation group (11) & $0.057^{***}$ & $0.108^{***}$ & $0.051^{***}$ & $0.377^{***}$ & $0.045^{***}$ & $0.318^{***}$ \\
 & (0.007) & (0.013) & (0.010) & (0.078) & (0.015) & (0.111) \\
Occupation group + industry group (11 + 33) & $0.078^{***}$ & $0.129^{***}$ & $0.050^{***}$ & $0.376^{***}$ & $0.050^{**}$ & $0.352^{***}$ \\
 & (0.010) & (0.017) & (0.012) & (0.087) & (0.020) & (0.129) \\
Cells ranked by pay and later growth (4 sets) & $0.056^{***}$ & $0.093^{***}$ & $0.036^{***}$ & $0.271^{***}$ & $0.033^{***}$ & $0.237^{***}$ \\
 & (0.006) & (0.012) & (0.009) & (0.067) & (0.011) & (0.087) \\
Occupation $\times$ industry cell (11 $\times$ 33) & $0.091^{***}$ & $0.146^{***}$ & $0.054^{***}$ & $0.405^{***}$ & $0.070^{***}$ & $0.493^{***}$ \\
 & (0.011) & (0.020) & (0.015) & (0.099) & (0.025) & (0.171) \\
Detailed occupation (occ1990dd) & $0.086^{***}$ & $0.165^{***}$ & $0.079^{***}$ & $0.591^{***}$ & $0.101^{***}$ & $0.713^{***}$ \\
 & (0.011) & (0.020) & (0.017) & (0.120) & (0.026) & (0.211) \\
\bottomrule
\end{tabular}}}
\par\smallskip
\begin{minipage}{\textwidth}\footnotesize \textit{Notes:} Short regression: the son's log wage on log parental income in each age band (25--34, 35--55), controls as in Table~\ref{tab:industrydecomp}. Long regression: adds fixed effects for the son's job at 25--34 (row), estimated separately in each age band and race sample, so they absorb any wage profile of where he started, linear or not. Accounted for = short $-$ long. On this sample the pooled coefficient rises by 0.132$^{***}$ and the White $-$ Black gap changes by 0.140$^{***}$ from 25--34 to 35--55. Ranked sets: cells (or industries) in the top quartile of men's employment by CPS wage level and by CPS wage growth after 35, and their intersection. Permuting the cell labels across sons gives shares of about 0.03. The son's early job also reflects his own skills, so the rows measure how much where a son stands at 25--34 predicts how far he pulls away, not a causal payoff. 93\% of sons observed at 35--55 have a job at 25--34. Standard errors clustered by 1968 family. $^{***}$ $p<0.01$, $^{**}$ $p<0.05$, $^{*}$ $p<0.1$. Source: \texttt{35.EarlyIndustryLatePayoff.py} (\texttt{output/early\_industry\_late\_payoff.csv}); table built by \texttt{28.PaperBTables.py}.\end{minipage}
\end{table}


\paragraph{Access to the jobs that build the gap.} Who gets these jobs differs by background and by race
(Table~\ref{tab:access}). Of sons from the top third of the parental income distribution, 20.2 percent start in cells
that are well paid with high later growth, against 4.7 percent from the bottom third; of White sons, 16.1 percent
against 5.7 percent of Black sons. The racial gap in access is not only a gap in parental income: at equal parental
income White sons are 4.2 percentage points more likely to start in these jobs (0.042$^{***}$), and 3.5 points
(0.035$^{*}$, $p=0.083$) holding the father's job and the son's education fixed as well. Parental income buys access more
steeply for White than for Black sons (a difference of 0.055$^{***}$ per log point), but not once education and the
father's job are held fixed ($-0.022$, s.e.\ 0.024). A father who held such a job raises his son's chance of holding one
by 6.4 percentage points (0.064$^{**}$; White sons 0.048$^{**}$, Black sons 0.008).\footnote{Son-level linear probability
models on the son's job at 25--34, clustered by 1968 family. For well-paid starting jobs of any kind, sons of rich White
parents are 24.7 points more likely than sons of rich Black parents to start in one (0.247$^{***}$;
\texttt{output/early\_industry\_late\_payoff.md}, Section D3).}

\begin{table}[!htbp]
\centering
\caption{Access to well-paid, fast-growing starting jobs at 25--34 (PSID, sons)}
\label{tab:access}
\resizebox{\ifdim\width>\textwidth\textwidth\else\width\fi}{!}{%
\resizebox*{!}{\ifdim\totalheight>0.85\textheight 0.85\textheight\else\totalheight\fi}{%
\begin{tabular}{lccccc}
\toprule
 & \begin{tabular}[b]{@{}c@{}}Well-paid,\\high later\\growth\\(a)\end{tabular} & \begin{tabular}[b]{@{}c@{}}Well-paid,\\high later\\growth\\(b)\end{tabular} & \begin{tabular}[b]{@{}c@{}}Well-paid,\\high later\\growth\\(c)\end{tabular} & \begin{tabular}[b]{@{}c@{}}Well-paid,\\high later\\growth\\(d)\end{tabular} & \begin{tabular}[b]{@{}c@{}}Well-paid,\\high later\\growth\\(e)\end{tabular} \\
 & (1) & (2) & (3) & (4) & (5) \\
\midrule
\multicolumn{6}{l}{\textit{Black and White sons (reference for terciles: Black, bottom tercile)}} \\
ln parental income & $0.104^{***}$ &  & $0.095^{***}$ &  & $0.002$ \\
 & (0.010) &  & (0.011) &  & (0.016) \\
White &  & $0.109^{***}$ & $0.042^{***}$ &  & $0.035^{*}$ \\
 &  & (0.014) & (0.016) &  & (0.020) \\
Black, middle tercile &  &  &  & $0.024$ &  \\
 &  &  &  & (0.022) &  \\
Black, top tercile &  &  &  & $0.125^{***}$ &  \\
 &  &  &  & (0.044) &  \\
White, bottom tercile &  &  &  & $0.048^{**}$ &  \\
 &  &  &  & (0.020) &  \\
White, middle tercile &  &  &  & $0.092^{***}$ &  \\
 &  &  &  & (0.017) &  \\
White, top tercile &  &  &  & $0.187^{***}$ &  \\
 &  &  &  & (0.017) &  \\
Rich White $-$ rich Black (top terciles) &  &  &  & $0.062$ &  \\
 &  &  &  & (0.045) &  \\
Father held a well-paid, high-later-growth job &  &  &  &  & $0.064^{**}$ \\
 &  &  &  &  & (0.025) \\
Years of education &  &  &  &  & $0.063^{***}$ \\
 &  &  &  &  & (0.005) \\
\addlinespace
\multicolumn{6}{l}{\textit{Black sons}} \\
ln parental income & $0.048^{***}$ &  & $0.048^{***}$ &  & $0.021$ \\
 & (0.009) &  & (0.009) &  & (0.018) \\
Middle tercile &  &  &  & $0.022^{*}$ &  \\
 &  &  &  & (0.013) &  \\
Top tercile &  &  &  & $0.104^{***}$ &  \\
 &  &  &  & (0.025) &  \\
Father held a well-paid, high-later-growth job &  &  &  &  & $0.008$ \\
 &  &  &  &  & (0.051) \\
Years of education &  &  &  &  & $0.052^{***}$ \\
 &  &  &  &  & (0.007) \\
\addlinespace
\multicolumn{6}{l}{\textit{White sons}} \\
ln parental income & $0.103^{***}$ &  & $0.103^{***}$ &  & $-0.001$ \\
 & (0.012) &  & (0.012) &  & (0.016) \\
Middle tercile &  &  &  & $0.050^{***}$ &  \\
 &  &  &  & (0.018) &  \\
Top tercile &  &  &  & $0.138^{***}$ &  \\
 &  &  &  & (0.019) &  \\
Father held a well-paid, high-later-growth job &  &  &  &  & $0.048^{**}$ \\
 &  &  &  &  & (0.023) \\
Years of education &  &  &  &  & $0.064^{***}$ \\
 &  &  &  &  & (0.005) \\
\bottomrule
\end{tabular}}}
\par\smallskip
\begin{minipage}{\textwidth}\footnotesize \textit{Continued on the next table; notes below its last part.}\end{minipage}
\end{table}
\begin{table}[!htbp]
\centering
\addtocounter{table}{-1}
\caption{Access to well-paid, fast-growing starting jobs at 25--34 (PSID, sons) (continued)}
\resizebox{\ifdim\width>\textwidth\textwidth\else\width\fi}{!}{%
\resizebox*{!}{\ifdim\totalheight>0.85\textheight 0.85\textheight\else\totalheight\fi}{%
\begin{tabular}{lccccc}
\toprule
 & \begin{tabular}[b]{@{}c@{}}Well-paid,\\high later\\growth\\(a)\end{tabular} & \begin{tabular}[b]{@{}c@{}}Well-paid,\\high later\\growth\\(b)\end{tabular} & \begin{tabular}[b]{@{}c@{}}Well-paid,\\high later\\growth\\(c)\end{tabular} & \begin{tabular}[b]{@{}c@{}}Well-paid,\\high later\\growth\\(d)\end{tabular} & \begin{tabular}[b]{@{}c@{}}Well-paid,\\high later\\growth\\(e)\end{tabular} \\
 & (1) & (2) & (3) & (4) & (5) \\
\midrule
\multicolumn{6}{l}{\textit{White $-$ Black}} \\
ln parental income & $0.055^{***}$ &  & $0.055^{***}$ &  & $-0.022$ \\
 & (0.015) &  & (0.015) &  & (0.024) \\
Middle tercile &  &  &  & $0.028$ &  \\
 &  &  &  & (0.022) &  \\
Top tercile &  &  &  & $0.034$ &  \\
 &  &  &  & (0.031) &  \\
Father held a well-paid, high-later-growth job &  &  &  &  & $0.040$ \\
 &  &  &  &  & (0.056) \\
Years of education &  &  &  &  & $0.012$ \\
 &  &  &  &  & (0.008) \\
\bottomrule
\end{tabular}}}
\par\smallskip
\begin{minipage}{\textwidth}\footnotesize \textit{Notes:} Son-level linear probability models. Outcome in columns (1)--(5): the son's job at 25--34 (modal occupation $\times$ industry cell) is in the top quartile of men's employment by both CPS wage level and CPS wage growth after 35 (mean 0.156); the same specifications for well-paid jobs of any kind are in \texttt{output/early\_industry\_late\_payoff.md}, Section D3. Specifications (a) log parental income; (b) White; (c) both; (d) race $\times$ parental-income tercile; (e) (c) plus the father's job and the son's years of education. Pooled: weighted; by race: unweighted. Standard errors clustered by 1968 family. $^{***}$ $p<0.01$, $^{**}$ $p<0.05$, $^{*}$ $p<0.1$. Source: \texttt{IGE/35.EarlyIndustryLatePayoff.py} (\texttt{output/early\_industry\_late\_payoff.csv}); table built by \texttt{IGE/28.PaperBTables.py}.\end{minipage}
\end{table}


\subsection{Wage growth with the same employer}

Using pairs of consecutive interviews one to two years apart at which the son stayed with the same employer, I regress
the annualised growth of his log wage on parental income (Table~\ref{tab:growthladder}). For White sons at 25--34 the
gradient is 0.021$^{***}$ (0.006): a White son whose parents had twice the income sees his wage grow about 1.4 percentage
points a year faster with the same employer, on mean growth of 5.7 percent. For Black sons it is 0.003 (0.008), and the
racial difference, 0.018$^{*}$ ($p=0.062$), is less precisely estimated than its components. The White gradient is not
measurement: over two intervals it is 0.019$^{***}$ and among strict stayers whose tenure covers both income years
0.013$^{**}$.\footnote{Controlling for the wage at the interview before the pair raises it to 0.036$^{***}$, so it is not
catch-up after a low start; over whole employer spells it is 0.019$^{***}$. \texttt{output/same\_employer\_growth\_checks.md},
Section B. Code: \texttt{IGE/36.SameEmployerGrowthChecks.py}.} About half of it is sorting into steep jobs: within the
cell of the first interview of the pair it is 0.011$^{*}$ ($p=0.095$), within the cell of the later interview 0.006
($p=0.34$), and within cells the racial difference is 0.006--0.007, not significant. Controlling for the son's years of
education lowers it to 0.008 ($p=0.17$). In the NLSY79, where the employer is observed directly in the job roster, the
White gradient is 0.008$^{*}$ and the Black gradient zero, and holding the AFQT and education fixed brings the White
gradient to 0.002: the part explained runs through education (0.0046$^{***}$), not measured ability (0.0014, $p=0.31$).
The precise statement is: for White sons the gradient is precisely estimated, half sorting and half growth within jobs,
and mostly schooling; for Black sons it is a precise zero; the racial difference in growth within the same job is not
established.\footnote{Annual growth between two interviews is too noisy to settle the part that does not run through
schooling. Cumulated over ten and twenty years of a son's own wage, the PSID gradient at equal schooling is
0.073$^{***}$ and 0.135$^{***}$ for White sons and about zero for Black sons with the 13--19 mean of parental income,
and about zero with a one-year measure, in the PSID as in the NLSY79 and NLSY97 (Section~\ref{sec:model},
Tables~\ref{tab:h4longdiff} and \ref{tab:h4measure}).}

\begin{table}[!htbp]
\centering
\caption{Wage growth with the same employer at 25--34: parental income, schooling and ability}
\label{tab:growthladder}
\resizebox{\ifdim\width>\textwidth\textwidth\else\width\fi}{!}{%
\resizebox*{!}{\ifdim\totalheight>0.85\textheight 0.85\textheight\else\totalheight\fi}{%
\begin{tabular}{lcccccccccc}
\toprule
 & \begin{tabular}[b]{@{}c@{}}PSID\\baseline\end{tabular} & \begin{tabular}[b]{@{}c@{}}PSID\\cell FE,\\first interview\end{tabular} & \begin{tabular}[b]{@{}c@{}}PSID\\cell FE,\\later interview\end{tabular} & \begin{tabular}[b]{@{}c@{}}PSID\\+ education\end{tabular} & \begin{tabular}[b]{@{}c@{}}PSID\\two-interval\\growth\end{tabular} & \begin{tabular}[b]{@{}c@{}}PSID\\strict\\stayers\end{tabular} & \begin{tabular}[b]{@{}c@{}}NLSY79\\baseline\end{tabular} & \begin{tabular}[b]{@{}c@{}}NLSY79\\+ AFQT\end{tabular} & \begin{tabular}[b]{@{}c@{}}NLSY79\\+ education\end{tabular} & \begin{tabular}[b]{@{}c@{}}NLSY79\\+ AFQT\\+ education\end{tabular} \\
 & (1) & (2) & (3) & (4) & (5) & (6) & (7) & (8) & (9) & (10) \\
\midrule
\multicolumn{11}{l}{\textit{White sons}} \\
Log parental income & $0.021^{***}$ & $0.011^{*}$ & $0.006$ & $0.008$ & $0.019^{***}$ & $0.013^{**}$ & $0.008^{*}$ & $0.004$ & $0.002$ & $0.002$ \\
 & (0.006) & (0.006) & (0.007) & (0.006) & (0.005) & (0.006) & (0.005) & (0.005) & (0.005) & (0.005) \\
Years of education &  &  &  & $0.009^{***}$ &  &  &  &  & $0.004^{***}$ & $0.004^{***}$ \\
 &  &  &  & (0.002) &  &  &  &  & (0.001) & (0.001) \\
AFQT (1 s.d.) &  &  &  &  &  &  &  & $0.009^{***}$ &  & $0.003$ \\
 &  &  &  &  &  &  &  & (0.003) &  & (0.003) \\
\addlinespace
\multicolumn{11}{l}{\textit{Black sons}} \\
Log parental income & $0.003$ & $0.004$ & $-0.001$ & $0.002$ & $-0.001$ & $0.010$ & $0.000$ & $0.000$ & $-0.002$ & $-0.001$ \\
 & (0.008) & (0.010) & (0.009) & (0.008) & (0.007) & (0.008) & (0.005) & (0.005) & (0.005) & (0.005) \\
Years of education &  &  &  & $0.002$ &  &  &  &  & $0.003$ & $0.004$ \\
 &  &  &  & (0.003) &  &  &  &  & (0.002) & (0.002) \\
AFQT (1 s.d.) &  &  &  &  &  &  &  & $0.003$ &  & $-0.003$ \\
 &  &  &  &  &  &  &  & (0.005) &  & (0.005) \\
\addlinespace
\multicolumn{11}{l}{\textit{White $-$ Black}} \\
Log parental income & $0.018^{*}$ & $0.006$ & $0.007$ & $0.006$ & $0.019^{**}$ & $0.003$ & $0.009$ & $0.004$ & $0.004$ & $0.003$ \\
 & (0.010) & (0.012) & (0.011) & (0.010) & (0.008) & (0.010) & (0.007) & (0.007) & (0.007) & (0.007) \\
Years of education &  &  &  & $0.006^{*}$ &  &  &  &  & $0.001$ & $0.000$ \\
 &  &  &  & (0.003) &  &  &  &  & (0.002) & (0.003) \\
AFQT (1 s.d.) &  &  &  &  &  &  &  & $0.006$ &  & $0.006$ \\
 &  &  &  &  &  &  &  & (0.006) &  & (0.006) \\
\bottomrule
\end{tabular}}}
\par\smallskip
\begin{minipage}{\textwidth}\footnotesize \textit{Notes:} Pairs of consecutive interviews one to two years apart at which the son is with the same employer (PSID: tenure covers the interval; NLSY79: the same employer in the job roster), sons aged 25--34 at the later interview. Outcome: annualised growth of the log hourly wage. One regression per column and race; controls in every column: quartic in age, year (round) fixed effects. Cell: 11 occupation groups $\times$ 33 industry groups. Two-interval growth: the same employer at three interviews. Strict stayers: tenure covers both income years. NLSY79 columns: common sample with AFQT and education known; AFQT standardised within birth year. White $-$ Black: difference of the separate regressions, joint influence functions, clustered by 1968 family (PSID) or 1979 household (NLSY79). $^{***}$ $p<0.01$, $^{**}$ $p<0.05$, $^{*}$ $p<0.1$. Source: \texttt{IGE/36.SameEmployerGrowthChecks.py} (\texttt{output/same\_employer\_growth\_checks.csv}).\end{minipage}
\end{table}


\subsection{Schooling: a return that rises with age, faster for White sons}\label{sec:schooling}

If schooling carries the gradient, the distinctive prediction is that its pay-off grows over the career.
\citet{bhuller2017life} show in Norwegian registers that education premiums rise steeply with age.
Table~\ref{tab:schooling} estimates the return to a year of schooling in the PSID by age band and race. For White sons the
return in hourly wages rises from 0.049 at 25--29 to 0.128 at 40--44, and in earnings from 0.045 to 0.143; for Black sons
it starts higher---0.076 in wages and 0.124 in earnings at 25--29---and rises much less in wages and not at all in
earnings (0.12--0.15 at every age). The White--Black difference in the rise is 0.032$^{***}$ (0.011) in wages and
0.055$^{***}$ (0.014) in earnings, and 0.013$^{*}$ (0.008) for the NLSY79 hourly rate of pay. Education accounts for a
large part of the divergence: adding years of schooling to the regression of the log wage on parental income by band
removes 0.057$^{***}$ of the pooled rise of wage persistence after 35 (a share of 0.44$^{***}$), and 0.069$^{***}$ of the
0.186$^{***}$ change in the racial gap (0.37$^{***}$); for earnings 0.090$^{***}$ of 0.203$^{***}$ (0.44$^{***}$), and in
the NLSY79 0.038$^{***}$ of 0.054$^{**}$ (0.72$^{**}$). Education is an outcome of the same background, so these are
shares of accounting, not causal mediation.\footnote{Code: \texttt{IGE/42.SchoolingAndJobGrowth.py}
(\texttt{output/schooling\_job\_growth.md}).}

\begin{table}[!htbp]
\centering
\caption{The return to schooling by age and race, and the part of wage persistence it carries}
\label{tab:schooling}
\resizebox{\ifdim\width>\textwidth\textwidth\else\width\fi}{!}{%
\resizebox*{!}{\ifdim\totalheight>0.85\textheight 0.85\textheight\else\totalheight\fi}{%
\begin{tabular}{lccccccc}
\toprule
 & \begin{tabular}[b]{@{}c@{}}PSID\\ln hourly wage\\on education\end{tabular} & \begin{tabular}[b]{@{}c@{}}PSID\\ln earnings\\on education\end{tabular} & \begin{tabular}[b]{@{}c@{}}NLSY79\\ln hourly wage\\on education\end{tabular} & \begin{tabular}[b]{@{}c@{}}PSID\\ln hourly wage\\on ln parental\\income\end{tabular} & \begin{tabular}[b]{@{}c@{}}PSID\\accounted for\\by education\end{tabular} & \begin{tabular}[b]{@{}c@{}}NLSY79\\ln hourly wage\\on ln parental\\income\end{tabular} & \begin{tabular}[b]{@{}c@{}}NLSY79\\accounted for\\by education\end{tabular} \\
 & (1) & (2) & (3) & (4) & (5) & (6) & (7) \\
\midrule
\multicolumn{8}{l}{\textit{White sons}} \\
Age 25--29 & $0.049^{***}$ & $0.045^{***}$ & $0.046^{***}$ & $0.220^{***}$ & $0.046^{***}$ & $0.222^{***}$ & $0.048^{***}$ \\
 & (0.005) & (0.007) & (0.004) & (0.018) & (0.008) & (0.018) & (0.007) \\
Age 30--34 & $0.090^{***}$ & $0.101^{***}$ & $0.072^{***}$ & $0.342^{***}$ & $0.101^{***}$ & $0.263^{***}$ & $0.082^{***}$ \\
 & (0.006) & (0.007) & (0.005) & (0.024) & (0.011) & (0.023) & (0.009) \\
Age 35--39 & $0.110^{***}$ & $0.122^{***}$ & $0.097^{***}$ & $0.384^{***}$ & $0.128^{***}$ & $0.319^{***}$ & $0.116^{***}$ \\
 & (0.007) & (0.008) & (0.005) & (0.030) & (0.014) & (0.025) & (0.012) \\
Age 40--44 & $0.128^{***}$ & $0.143^{***}$ & $0.120^{***}$ & $0.440^{***}$ & $0.153^{***}$ & $0.338^{***}$ & $0.149^{***}$ \\
 & (0.008) & (0.009) & (0.005) & (0.037) & (0.017) & (0.028) & (0.014) \\
Age 45--49 & $0.125^{***}$ & $0.130^{***}$ & $0.122^{***}$ & $0.465^{***}$ & $0.140^{***}$ & $0.348^{***}$ & $0.153^{***}$ \\
 & (0.010) & (0.011) & (0.006) & (0.050) & (0.020) & (0.031) & (0.015) \\
Age 50--55 & $0.128^{***}$ & $0.128^{***}$ & $0.123^{***}$ & $0.455^{***}$ & $0.162^{***}$ & $0.367^{***}$ & $0.146^{***}$ \\
 & (0.012) & (0.014) & (0.007) & (0.053) & (0.026) & (0.035) & (0.016) \\
\textbf{Change, 35+ minus below 35} & $0.053^{***}$ & $0.057^{***}$ & $0.056^{***}$ & $0.155^{***}$ & $0.072^{***}$ & $0.100^{***}$ & $0.076^{***}$ \\
 & (0.007) & (0.008) & (0.004) & (0.030) & (0.013) & (0.020) & (0.009) \\
\addlinespace
\multicolumn{8}{l}{\textit{Black sons}} \\
Age 25--29 & $0.076^{***}$ & $0.124^{***}$ & $0.068^{***}$ & $0.174^{***}$ & $0.037^{***}$ & $0.146^{***}$ & $0.047^{***}$ \\
 & (0.007) & (0.010) & (0.006) & (0.020) & (0.007) & (0.018) & (0.007) \\
Age 30--34 & $0.101^{***}$ & $0.151^{***}$ & $0.086^{***}$ & $0.199^{***}$ & $0.058^{***}$ & $0.178^{***}$ & $0.060^{***}$ \\
 & (0.008) & (0.010) & (0.008) & (0.023) & (0.010) & (0.020) & (0.010) \\
Age 35--39 & $0.106^{***}$ & $0.144^{***}$ & $0.105^{***}$ & $0.187^{***}$ & $0.062^{***}$ & $0.165^{***}$ & $0.074^{***}$ \\
 & (0.009) & (0.012) & (0.008) & (0.030) & (0.010) & (0.025) & (0.013) \\
Age 40--44 & $0.096^{***}$ & $0.127^{***}$ & $0.115^{***}$ & $0.140^{***}$ & $0.042^{***}$ & $0.225^{***}$ & $0.090^{***}$ \\
 & (0.011) & (0.014) & (0.010) & (0.034) & (0.011) & (0.027) & (0.016) \\
Age 45--49 & $0.118^{***}$ & $0.146^{***}$ & $0.127^{***}$ & $0.145^{***}$ & $0.057^{***}$ & $0.206^{***}$ & $0.095^{***}$ \\
 & (0.012) & (0.016) & (0.010) & (0.042) & (0.016) & (0.029) & (0.017) \\
Age 50--55 & $0.121^{***}$ & $0.143^{***}$ & $0.132^{***}$ & $0.150^{***}$ & $0.040^{**}$ & $0.238^{***}$ & $0.105^{***}$ \\
 & (0.016) & (0.021) & (0.010) & (0.057) & (0.019) & (0.032) & (0.017) \\
\textbf{Change, 35+ minus below 35} & $0.021^{**}$ & $0.002$ & $0.043^{***}$ & $-0.030$ & $0.003$ & $0.047^{***}$ & $0.038^{***}$ \\
 & (0.009) & (0.011) & (0.006) & (0.029) & (0.009) & (0.018) & (0.009) \\
\addlinespace
\multicolumn{8}{l}{\textit{White $-$ Black}} \\
Age 25--29 & $-0.027^{***}$ & $-0.079^{***}$ & $-0.022^{***}$ & $0.046^{*}$ & $0.009$ & $0.076^{***}$ & $0.000$ \\
 & (0.009) & (0.012) & (0.007) & (0.027) & (0.011) & (0.026) & (0.010) \\
Age 30--34 & $-0.011$ & $-0.049^{***}$ & $-0.013$ & $0.143^{***}$ & $0.043^{***}$ & $0.085^{***}$ & $0.023^{*}$ \\
 & (0.010) & (0.013) & (0.009) & (0.033) & (0.015) & (0.030) & (0.013) \\
Age 35--39 & $0.005$ & $-0.022$ & $-0.008$ & $0.197^{***}$ & $0.065^{***}$ & $0.154^{***}$ & $0.041^{**}$ \\
 & (0.012) & (0.015) & (0.010) & (0.043) & (0.017) & (0.035) & (0.017) \\
Age 40--44 & $0.032^{**}$ & $0.016$ & $0.004$ & $0.299^{***}$ & $0.111^{***}$ & $0.113^{***}$ & $0.060^{***}$ \\
 & (0.014) & (0.017) & (0.012) & (0.051) & (0.020) & (0.039) & (0.021) \\
Age 45--49 & $0.007$ & $-0.015$ & $-0.004$ & $0.320^{***}$ & $0.083^{***}$ & $0.142^{***}$ & $0.059^{***}$ \\
 & (0.016) & (0.020) & (0.011) & (0.066) & (0.025) & (0.043) & (0.022) \\
Age 50--55 & $0.008$ & $-0.015$ & $-0.009$ & $0.305^{***}$ & $0.122^{***}$ & $0.129^{***}$ & $0.041^{*}$ \\
 & (0.020) & (0.025) & (0.013) & (0.078) & (0.032) & (0.047) & (0.024) \\
\textbf{Change, 35+ minus below 35} & $0.032^{***}$ & $0.055^{***}$ & $0.013^{*}$ & $0.186^{***}$ & $0.069^{***}$ & $0.054^{**}$ & $0.038^{***}$ \\
 & (0.011) & (0.014) & (0.008) & (0.041) & (0.016) & (0.027) & (0.013) \\
\bottomrule
\end{tabular}}}
\par\smallskip
\begin{minipage}{\textwidth}\footnotesize \textit{Notes:} Son-years at 25--55 with a valid wage, years of education and parental income; one regression per column and age band, controls: quartic in age and year (round) fixed effects. Columns (1)--(3): coefficient on years of completed education. Columns (4) and (6): coefficient on log parental income. Columns (5) and (7): the part of (4) and (6) that disappears when years of education are added (the return to education times the education gradient in parental income). Change row: mean over the bands from 35 minus mean over the bands below 35. White $-$ Black: difference of the separate regressions, standard errors from joint influence functions, clustered by 1968 family (PSID) or 1979 household (NLSY79). $^{***}$ $p<0.01$, $^{**}$ $p<0.05$, $^{*}$ $p<0.1$. Source: \texttt{42.SchoolingAndJobGrowth.py} (\texttt{output/schooling\_job\_growth.csv}).\end{minipage}
\end{table}


\paragraph{Within jobs, net of the calendar, holding ability fixed.} A steeper White return could reflect access alone:
if good jobs accumulate human capital faster and educated White sons get more of them, the return to schooling rises
faster with age for White sons with no racial difference in what schooling does on a given job. The test holds the job
fixed (Table~\ref{tab:withincell}). On son-years with an early job, the White--Black difference in the rise of the
earnings return is 0.058$^{***}$ (0.014) without cell controls, 0.061$^{***}$ (0.018) with fixed effects for the son's
cell at 25--34 and 0.066$^{***}$ (0.016) with fixed effects for his current cell; in wages 0.031$^{***}$, 0.029$^{*}$ and
0.036$^{***}$. Holding the early job fixed halves the rise of the earnings return in the pooled sample (from 0.050$^{***}$ to
0.025$^{***}$), so half of the rising return is the sorting of the educated into steep jobs; but it lowers the rise for
both races alike, and the racial difference does not move (changes of $-0.002$ to 0.009, s.e.\ 0.011--0.016).
The rise of the college premium over 1980--2000 could make a return rise with age that is really a rise with calendar
time; allowing the return to trend with the calendar leaves the racial difference in the rise at 0.063$^{**}$ (0.027)
for earnings and 0.035 ($p=0.12$) for wages, and within the cohorts born 1955--69 it is 0.051$^{***}$ and 0.029$^{*}$. In
the NLSY79, a single cohort, the difference is 0.017$^{*}$ within the current cell and 0.021$^{**}$ (0.009) when the AFQT
is held fixed. The Black sons' higher return at 25--29 also survives the job (0.061$^{***}$ in earnings within the early
cell) in the PSID, but not in the NLSY79 (0.008).\footnote{Code: \texttt{IGE/45.WithinCellIdentification.py}
(\texttt{output/within\_cell\_identification.md}); NLSY79 AFQT results in Appendix Table~\ref{tab:afqtgrowth}.}

The age pattern of the returns says which mechanism this is. Under employer learning the coefficient on schooling, which
employers see, falls with experience, and the coefficient on ability, which they learn, rises
\citep{altonji2001employer, lange2007speed, arcidiacono2010beyond}. In the NLSY79 both rise: with the AFQT and schooling
in the same regression, the schooling coefficient rises by 0.049$^{***}$ from the bands below 35 to the bands above, and
the AFQT coefficient by 0.036$^{***}$, and the AFQT rise does not differ by race ($-0.031$, s.e.\ 0.025). A rising return
to schooling is the complementarity of schooling with human capital accumulated on the job \citep{bhuller2017life},
and it is this complementarity that is weaker for Black sons, within the same job and at equal measured ability.
The return to experience itself does not differ by race at equal AFQT: in the within-decile regression of
\citet{rauh2023black}, run on the NLSY79 with its supplemental Black sample, hourly earnings grow with experience at the
same rate for Black and White men (Black $\times$ experience 0.003 in logs, s.e.\ 0.003, $p = 0.30$; 1,184 Black men).
Their model holds accumulation race-neutral given the AFQT and has no schooling; the racial difference here is in how
growth varies with schooling.\footnote{Hourly earnings (2018 dollars, and logs) on Black, years of accumulated
experience, their product and region fixed effects by AFQT decile, men aged 23--54 not in school, at least 500 hours and
25 weeks in the year; standard errors clustered by person. On the cross-sectional sample (285 Black men) the regression
reproduces their pattern: no difference in deciles 1--5, and negative Black $\times$ experience in deciles 9 and 10
($-1.03^{**}$ and $-1.93^{***}$ dollars a year, $p = 0.034$ and $p < 0.001$), cells that hold 4 and 3 Black men. With
the supplemental sample no decile differs and the slope pooled within deciles is 0.028 dollars a year (s.e.\ 0.071,
$p = 0.70$). Code: \texttt{IGE/70.RauhReplication.py}; output \texttt{output/rauh\_replication.md}.}
The same-employer evidence points the same way but is weaker: a year of schooling raises same-employer wage growth at
25--34 by 0.010$^{***}$ for White and 0.003 for Black sons (difference 0.007$^{**}$, $p=0.030$), and within the cell of
the first interview the difference is 0.004 ($p=0.37$); in the NLSY79 it is 0.002 ($p=0.47$) and 0.005$^{*}$ within the
cell.

\begin{table}[!htbp]
\centering
\caption{Does the racial difference in the pay-off of schooling survive within jobs?}
\label{tab:withincell}
\resizebox{\ifdim\width>\textwidth\textwidth\else\width\fi}{!}{%
\resizebox*{!}{\ifdim\totalheight>0.85\textheight 0.85\textheight\else\totalheight\fi}{%
\begin{tabular}{lccccc}
\toprule
 & \begin{tabular}[b]{@{}c@{}}Baseline\end{tabular} & \begin{tabular}[b]{@{}c@{}}Early-cell\\fixed effects\end{tabular} & \begin{tabular}[b]{@{}c@{}}Current-cell\\fixed effects\end{tabular} & \begin{tabular}[b]{@{}c@{}}Schooling $\times$\\calendar trend\end{tabular} & \begin{tabular}[b]{@{}c@{}}Born\\1955--69\end{tabular} \\
 & (1) & (2) & (3) & (4) & (5) \\
\midrule
\multicolumn{6}{l}{\textit{PSID, log annual earnings}} \\
White: rise of the return, 35+ minus below 35 & $0.057^{***}$ & $0.035^{***}$ & $0.032^{***}$ & $0.041^{***}$ & $0.064^{***}$ \\
 & (0.008) & (0.010) & (0.009) & (0.015) & (0.010) \\
Black: rise of the return, 35+ minus below 35 & $-0.001$ & $-0.025^{*}$ & $-0.034^{***}$ & $-0.022$ & $0.013$ \\
 & (0.012) & (0.015) & (0.013) & (0.023) & (0.016) \\
White $-$ Black: rise of the return & $0.058^{***}$ & $0.061^{***}$ & $0.066^{***}$ & $0.063^{**}$ & $0.051^{***}$ \\
 & (0.014) & (0.018) & (0.016) & (0.027) & (0.019) \\
White $-$ Black: return at 25--29 & $-0.077^{***}$ & $-0.061^{***}$ & $-0.066^{***}$ & $-0.079^{***}$ & $-0.057^{***}$ \\
 & (0.012) & (0.015) & (0.012) & (0.012) & (0.019) \\
\addlinespace
\multicolumn{6}{l}{\textit{PSID, log hourly wage}} \\
White: rise of the return, 35+ minus below 35 & $0.052^{***}$ & $0.032^{***}$ & $0.028^{***}$ & $0.038^{***}$ & $0.060^{***}$ \\
 & (0.007) & (0.009) & (0.008) & (0.014) & (0.009) \\
Black: rise of the return, 35+ minus below 35 & $0.021^{**}$ & $0.003$ & $-0.008$ & $0.003$ & $0.031^{**}$ \\
 & (0.009) & (0.012) & (0.010) & (0.018) & (0.013) \\
White $-$ Black: rise of the return & $0.031^{***}$ & $0.029^{*}$ & $0.036^{***}$ & $0.035$ & $0.029^{*}$ \\
 & (0.011) & (0.015) & (0.013) & (0.023) & (0.016) \\
White $-$ Black: return at 25--29 & $-0.028^{***}$ & $-0.023^{**}$ & $-0.030^{***}$ & $-0.028^{***}$ & $-0.016$ \\
 & (0.009) & (0.010) & (0.009) & (0.009) & (0.014) \\
\addlinespace
\multicolumn{6}{l}{\textit{NLSY79, log hourly rate of pay}} \\
White: rise of the return, 35+ minus below 35 & $0.056^{***}$ & $0.044^{***}$ & $0.046^{***}$ &  &  \\
 & (0.004) & (0.006) & (0.005) &  &  \\
Black: rise of the return, 35+ minus below 35 & $0.044^{***}$ & $0.027^{***}$ & $0.029^{***}$ &  &  \\
 & (0.006) & (0.009) & (0.007) &  &  \\
White $-$ Black: rise of the return & $0.012$ & $0.017$ & $0.017^{*}$ &  &  \\
 & (0.008) & (0.010) & (0.009) &  &  \\
White $-$ Black: return at 25--29 & $-0.022^{***}$ & $-0.008$ & $-0.004$ &  &  \\
 & (0.007) & (0.008) & (0.007) &  &  \\
\bottomrule
\end{tabular}}}
\par\smallskip
\begin{minipage}{\textwidth}\footnotesize \textit{Notes:} Coefficient on years of schooling in a regression of the log outcome by age band (25--29 to 50--55) and race; rows give its rise (mean over the bands from 35 minus mean below 35) and, for the race difference, its value at 25--29. Controls: quartic in age, year (round) fixed effects. Early cell: the son's modal occupation group $\times$ industry group (11 $\times$ 33) at 25--34; current cell: the same at each observation. Column (4) adds schooling $\times$ (year $-$ 1995)/10, so the return is evaluated in 1995 and the rise is net of the calendar trend of the return; column (5) keeps sons born 1955--69. The NLSY79 is one cohort, so (4)--(5) are not defined there. All columns use the same son-years (early cell known). Standard errors clustered by 1968 family (PSID) or 1979 household (NLSY79). $^{***}$ $p<0.01$, $^{**}$ $p<0.05$, $^{*}$ $p<0.1$. Source: \texttt{45.WithinCellIdentification.py} (\texttt{output/within\_cell\_identification.csv}).\end{minipage}
\end{table}


\subsection{Wage growth within and between jobs}

Wage growth comes from staying, from moving and from losing jobs \citep{topel1992job, altonji2013modeling}.
Table~\ref{tab:jobmoves} splits the parental-income gradient of annual wage growth into these three parts. For White
sons at 25--34 the gradient of 0.020$^{***}$ is almost all within jobs: 0.017$^{***}$ comes from years with the same
employer, 0.001 from moves and 0.003 from displacements. For Black sons the within-job part is 0.005, and moves
contribute $-0.008^{**}$: Black sons of richer parents gain less, not more, when they change employers. Of the racial
difference in the growth gradient at 25--34, 0.021$^{***}$, about three fifths is within jobs (0.012$^{*}$) and two fifths
at moves (0.009$^{**}$); the racial difference in the gain at a move is 0.069$^{**}$ (0.028). Family background buys the
same protection from job loss in both races: the probability of displacement falls with parental income by 0.037$^{***}$
for White and 0.038$^{***}$ for Black sons.

\begin{table}[!htbp]
\centering
\caption{Wage growth within and between jobs by parental income (PSID, sons)}
\label{tab:jobmoves}
\resizebox{\ifdim\width>\textwidth\textwidth\else\width\fi}{!}{%
\resizebox*{!}{\ifdim\totalheight>0.85\textheight 0.85\textheight\else\totalheight\fi}{%
\begin{tabular}{lcccccccc}
\toprule
 & \begin{tabular}[b]{@{}c@{}}Wage growth,\\all pairs\end{tabular} & \begin{tabular}[b]{@{}c@{}}Contribution:\\stayed\end{tabular} & \begin{tabular}[b]{@{}c@{}}Contribution:\\moved\end{tabular} & \begin{tabular}[b]{@{}c@{}}Contribution:\\displaced\end{tabular} & \begin{tabular}[b]{@{}c@{}}Stayed\end{tabular} & \begin{tabular}[b]{@{}c@{}}Displaced\end{tabular} & \begin{tabular}[b]{@{}c@{}}Growth\\if stayed\end{tabular} & \begin{tabular}[b]{@{}c@{}}Growth\\if moved\end{tabular} \\
 & (1) & (2) & (3) & (4) & (5) & (6) & (7) & (8) \\
\midrule
\multicolumn{9}{l}{\textit{Sons aged 25--34}} \\
White: log parental income & $0.020^{***}$ & $0.017^{***}$ & $0.001$ & $0.003$ & $0.035^{***}$ & $-0.037^{***}$ & $0.020^{***}$ & $0.005$ \\
 & (0.005) & (0.004) & (0.002) & (0.002) & (0.012) & (0.009) & (0.006) & (0.017) \\
Black: log parental income & $-0.001$ & $0.005$ & $-0.008^{**}$ & $0.002$ & $0.035^{**}$ & $-0.038^{***}$ & $0.004$ & $-0.064^{***}$ \\
 & (0.006) & (0.005) & (0.003) & (0.004) & (0.016) & (0.014) & (0.008) & (0.023) \\
White $-$ Black: log parental income & $0.021^{***}$ & $0.012^{*}$ & $0.009^{**}$ & $0.000$ & $0.000$ & $0.001$ & $0.016$ & $0.069^{**}$ \\
 & (0.008) & (0.007) & (0.004) & (0.005) & (0.020) & (0.017) & (0.010) & (0.028) \\
\addlinespace
\multicolumn{9}{l}{\textit{Sons aged 35--55}} \\
White: log parental income & $0.007^{**}$ & $0.003$ & $0.004^{**}$ & $0.000$ & $0.033^{***}$ & $-0.021^{**}$ & $0.003$ & $0.046^{***}$ \\
 & (0.003) & (0.003) & (0.002) & (0.001) & (0.011) & (0.008) & (0.004) & (0.017) \\
Black: log parental income & $-0.002$ & $0.000$ & $-0.002$ & $-0.001$ & $-0.005$ & $0.000$ & $0.000$ & $-0.013$ \\
 & (0.004) & (0.003) & (0.002) & (0.003) & (0.015) & (0.011) & (0.004) & (0.020) \\
White $-$ Black: log parental income & $0.009^{*}$ & $0.003$ & $0.006^{**}$ & $0.000$ & $0.039^{**}$ & $-0.021$ & $0.002$ & $0.059^{**}$ \\
 & (0.005) & (0.004) & (0.003) & (0.003) & (0.019) & (0.014) & (0.005) & (0.026) \\
\bottomrule
\end{tabular}}}
\par\smallskip
\begin{minipage}{\textwidth}\footnotesize \textit{Notes:} Consecutive interviews one to two years apart with a valid wage, a mapped occupation and tenure (waves 1985--2023). Between the two wage years the son stayed with the employer, moved without displacement, or was displaced (plant closure, layoff or firing). Columns (2)--(4) are the wage growth times an indicator of the event and add up to column (1) on the same pairs (column (3) and (4) of the moved and displaced contributions); (5)--(6) linear probabilities; (7)--(8) growth conditional on the event. Each cell: coefficient on log parental income; controls: quartic in age, year fixed effects; standard errors clustered by 1968 family. $^{***}$ $p<0.01$, $^{**}$ $p<0.05$, $^{*}$ $p<0.1$. Source: \texttt{IGE/42.SchoolingAndJobGrowth.py} (\texttt{output/schooling\_job\_growth.csv}).\end{minipage}
\end{table}


\subsection{What does not carry the divergence}

Entry into supervisory and managerial occupations, the closest observable analogue to the promotion delays documented
by \citet{hegde2026racial}, rises with parental income after 35 for White and not for Black sons (a racial difference of
0.024$^{***}$ at 35--55), but the change of that difference relative to 25--34 is 0.008 (0.011), and self-reported
position changes and promotions show no Black deficit. Job loss is frequent and costly---displacement lowers the son's log
wage two to three years later by 0.131$^{***}$, of which 0.028$^{***}$ is a lower occupational score---but the
occupational loss is not larger for Black sons, and shocks do not account for the divergence (Section~\ref{sec:contains}).
These results, the industry decompositions and the NLSY79 replication of the occupational split are in
Appendix~\ref{app:careers}.

\paragraph{Summary.} The divergence is built at 25--34. Sons of richer White parents start in jobs whose wages grow,
to which access is unequal by parental income and by race at equal parental income, and their wages grow faster in
them; schooling carries most of both, and its return rises with age faster for White than for Black sons within the same
job. Family background protects both races equally from displacement. The advantage built by the mid-thirties is what
the IGE measures from 35 on.
